\documentclass[
  a4paper,
  12pt,
  oneside,
  parskip=half,
  listof=totoc,
  bibliography=totoc,
  DIV=12,
  numbers=noenddot
  ]{scrartcl}

% ========== Grundlegende Kodierung & Schriften ==========
\usepackage[utf8]{inputenc}
\usepackage[T1]{fontenc}           % Für LuaLaTeX: ermöglicht Systemfonts und Unicode

% ========== Seitenlayout ==========
\usepackage[
  a4paper,
  left=2.5cm,       % Linker Rand = 2.5cm
  right=2.5cm,    % Rechter Rand = 2.5cm
  top=2.5cm,      % Oberer Rand = 2.5cm
  bottom=3cm,   % Unterer Rand = 3cm
  footskip=1.25cm,   % Abstand Fußzeile
  %bindingoffset=5mm % Bindungskorrektur
]{geometry}

\usepackage[onehalfspacing]{setspace}   % 1.5-facher Zeilenabstand
%\renewcommand*\familydefault{\sfdefault} % Serifenlose Schriftart
\usepackage{lmodern}               % Latin Modern Schrift (skalierbare Vektorfonts)
\usepackage[babel]{microtype}             % Verbesserte Textformatierung

% ========== Kopf- und Fußzeilen für scrreprt ==========
%\usepackage[headsepline]{scrlayer-scrpage}
%\automark[section]{chapter}   % inner header = section, outer = chapter
%\\textbf{}setkomafont{pagehead}{\normalfont}
%\pagestyle{scrheadings}
%\clearpairofpagestyles
%\ihead{\leftmark}    % inner header = Kapitel
%\ohead{\pagemark}    % outer header = Seitenzahl


% ========== Bibliographie ==========
\usepackage[
  backend=biber,       % Biber als Backend
  style=numeric-comp,       % Physik-Stil (AIP/APS Standard)
  sorting=none,        % Sortiere nach Reihenfolge des Auftretens im Text
  maxnames=3,         % Maximale Namen vor "et al."
  minnames=2          % Minimale Namen vor "et al."
  %doi=false,
  %url=false, 
  %isbn=false
  ]{biblatex}
\DeclareFieldFormat{titlecase}{#1} % Titel in Groß- und Kleinschreibung
\AtEveryBibitem{\clearfield{note}} %Weglassen von note
\AtEveryBibitem{\clearfield{pagetotal}} %Weglassen von Seitenzahlen
\DeclareFieldFormat[article]{journaltitle}{\textit{#1}} %Journal kursiv
\DeclareNameAlias{author}{family-given} %Reihenfolge Nachmane, Vorname tauschen
\DefineBibliographyStrings{english}{andothers = {\textit{et al\adddot}}} %et al. kursiv
\DeclareFieldAlias{journaltitle}{shortjournal} % Use shortjournal if available for abbreviated journal titles

\addbibresource{BA.bib} %Bibliographie-Datei einbinden


% ========== Sprach- & Zeitoptionen ==========
\usepackage[english]{babel}
\usepackage{csquotes}            % Für korrekte Anführungszeichen (automatisch an Sprache angepasst)

% ========== Naturwissenschaftliche Pakete ==========
\usepackage{amsmath}
\usepackage{amssymb}
\usepackage{mathtools}
\usepackage{bm}
\usepackage{fixmath}
%\usepackage{unicode-math}
\usepackage{physics}
%\usepackage{mhchem}            % Chemische Formeln
\usepackage[locale=DE]{siunitx} % SI-Einheiten mit deutschen Standards
\sisetup{
  output-decimal-marker = {,}, % Komma als Dezimaltrennzeichen
  per-mode = fraction,         % Einheiten als Bruch
  group-separator = {.},       % Tausendertrenner (Punkt)
  range-units = single         % Einheiten bei Bereichen: \SI{5-10}{\meter}
}
\DeclareSIUnit\pixel{px}


% ========== Grafikpakete ==========
\usepackage{graphicx}
\usepackage{pdfpages}               % PDF Impoort
%\usepackage{svg}
\usepackage{float}                  % Präzise Positionierung von Abbildungen
\usepackage{wrapfig}                % Floating Text
\usepackage{tikz}                   % Erstellung von Diagrammen und Grafiken
\usepackage[font=footnotesize,labelfont=bf]{caption}
\usepackage{mdframed}
\usepackage{pgfplots}
%\graphicspath{{./graphics}}

% ========== Tabellen ==========
\usepackage{tabularx}
\usepackage{array}                 % Erweiterte Tabellenfunktionen
\usepackage{booktabs}              % Professionelle Tabellenlinien
\usepackage{longtable}             % Mehrseitige Tabellen
\usepackage{multirow}              % Zellen über mehrere Zeilen
\setlength{\arrayrulewidth}{0.4mm} % Dicke der Tabellenlinien

% ========== Referenzen und Zitationen ==========
\usepackage{xurl}
\usepackage[colorlinks=true, urlcolor=black, linkcolor=black, anchorcolor=black, citecolor=blue]{hyperref}             
\usepackage[labelfont={bf}, labelsep=colon]{caption} % Anpassung von Bildunterschriften
\DeclareCaptionLabelSeparator{colon}{\space|\space} % Figurename | bei Abbildungen
\usepackage{subcaption}             % Unterbilder (a), (b), etc.
\usepackage{cleveref}
\usepackage{subfiles}

% ========== Benutzerdefinierte Befehle ==========
% Kurzbefehle für häufig verwendete Notationen
\newcommand{\di}{\partial}           % Partielle Ableitung (kurz)
\newcommand{\e}{\mathrm{e}}          % Euler's e (aufrecht)
\newcommand{\im}{\mathrm{i}}         % Imaginäre Einheit (aufrecht)
\newcommand{\dif}{\mathop{}\!\mathrm{d}} % Differential (richtige Abstände)
\newcommand{\HRule}{\rule{\linewidth}{0.5mm}} % Horizontale Linie (für Titelseite)
\newcommand{\ShouldBe}{\stackrel{\mathclap{\normalfont\mbox{!}}}{=}}
\newcommand{\bunit}[1]{\,\unit{#1}}

\begin{document}
	
\thispagestyle{empty}

\begin{center}
	\setlength{\tabcolsep}{0pt}
	\begin{tabular}{>{\raggedleft}m{2.5cm}>{\centering}m{\dimexpr\textwidth - 5cm\relax}>{\raggedright}m{2.5cm}}
		\includegraphics[width=\linewidth]{Media/Logos/Universität_Regensburg_logo.png}%
		&%
		
		&%
		\includegraphics[width=\linewidth]{Media/Logos/GSI-Logo-rgb.png} %
	\end{tabular}
\end{center}


\qquad

\setstretch{1.5}

\begin{center}
	
BACHELOR'S THESIS
	
\textbf{\huge Study of jet diffusion wakes in Monte Carlo simulations and in Pb--Pb collisions with ALICE}

SUBMITTED IN FULFILLMENT OF THE REQUIREMENTS FOR THE DEGREE OF\linebreak
BACHELOR OF SCIENCE\linebreak
IN PHYSICS

\HRule

by Nicolas Wirth

\qquad

CONDUCTED AT\linebreak
THE INSTITUTE OF THEORETICAL PHYSICS, UNIVERSITY OF REGENSBURG\linebreak
GSI HELMHOLTZ CENTER FOR HEAVY ION RESEARCH, DARMSTADT

\qquad

SUPERVISORS

\setstretch{1}

\begin{minipage}[t]{0.49\textwidth}
	\raggedright
	Prof. Dr. Sara Collins\linebreak
	Institute of Theoretical Physics\linebreak
	Faculty of Physics\linebreak
	University of Regensburg\linebreak
	
\end{minipage}
\begin{minipage}[t]{0.49\textwidth}
	\raggedleft
	Dr. Alexander Schmah\linebreak
	ALICE department\linebreak
	GSI\linebreak
	Darmstadt
\end{minipage}

\qquad

Regensburg, \today

\end{center}


\newpage

\thispagestyle{empty}
\mbox{}

\newpage

\setstretch{1}

\section*{}

\setcounter{page}{1}

\begin{center}
	\textbf{ABSTRACT}
\end{center}

In the early moments of our universe, matter was hot and dense, forming a state of matter referred to as quark-gluon plasma (QGP), a relativistic fluid in which partons are quasi-free. Ultra-relativistic heavy-ion collisions recreate droplets of the QGP, allowing studies of the QGP and emergent phenomena of quantum chromodynamics. Hard parton interactions between nuclei produce high-momentum, directed particle sprays, so called jets, that can penetrate the QGP. As they traverse the QGP medium, they lose energy and modify the QGP medium likewise. A study of these interactions reveals properties of the QGP. An interesting effect, the so called diffusion wake, was predicted by relativistic QGP fluid dynamics. A depletion of soft hadrons in the opposite direction relative to the jet direction is predicted. This depletion has been measured by a few groups experimentally in different ways. In this work, the diffusion wake is measured using the correlation of soft hadrons with a dijet/dihadron trigger pair in ALICE Run 3 Pb--Pb data. Furthermore, a simulation of ALICE Pb--Pb data is used to identify possible biases this analysis is subject to.

\newpage

\mbox{}

\newpage

\tableofcontents

\newpage

\mbox{}

\newpage

\newpage

\subfile{SECTION_QGP_physics_at_colliders.tex}

\pagebreak

\subfile{SECTION_ALICE_at_the_lhc.tex}

\pagebreak

\subfile{SECTION_Ana_of_heavy_ion_data.tex}

\pagebreak

\subfile{SECTION_measurements_in_ALICE_data.tex}

\pagebreak

\subfile{SECTION_event_simulation.tex}

%\pagebreak

%\subfile{SECTION_simulation_of_jet_shift.tex}

\pagebreak

\subfile{SECTION_summary.tex}

\pagebreak

\subfile{SECTION_appendix.tex}

\pagebreak

\printbibliography{}

\pagebreak

%\listoffigures

\pagebreak

\section*{Acknowledgements}


I would like to thank a few people, but foremost my excellent supervisor at GSI, Dr. Alexander Mark Schmah, for his daily lectures on heavy ion physics and critical thinking. I would also like to appreciate his open-mindedness and his ideas, which led me to experiments that turned out to be extraordinarily fruitful. I would also like to thank Nicola Wilson, who supported me especially at the beginning and at the end of my thesis. I would probably still be trying to obtain Large Hadron Collider data without her valuable knowledge. Prof. Dr. Peter M. Jacobs has contributed greatly to what I understand of heavy-ion collisions now compared to when I started. I would like to thank him for his effort on supporting and advising me. Another person crucial to this thesis is my supervisor from the University of Regensburg, Prof. Dr. Sara Collins, whom I would like to thank. For her immediate support on writing a thesis in experimental physics at her chair of theoretical physics. This openness shows great character and is admirable and indispensable for the presented work. I am especially grateful for the loving support of my mother who has made all this possible.


\pagebreak

\section*{Declaration of Authorship}

Ich, Nicolas Wirth, geb. am 27. 05. 2002, habe diese Arbeit selbstständig verfasst und keine anderen als die angegebenen Quellen und Hilfsmittel dafür benutzt. Die Arbeit ist nicht bereits bei einer anderen Hochschule zur Erlangung eines akademischen Grades eingereicht worden. Ich bin über wissenschaftlich korrektes Arbeiten aufgeklärt worden und habe von den in \S 24 Abs. 6 vorgesehenen Rechtsfolgen Kenntnis genommen. Die vorgelegten Druckexemplare und die vorgelegte elektronische Version der Arbeit sind identisch.

\qquad

\qquad

\qquad

\qquad


------------------------------------------

Regensburg, \today

\end{document}